\input{preamble}

\title{Section 4.2 --- The Logarithm Function}

\begin{document}
\maketitle

\section{Translating Between Exponentials and Logs}
Write each equation in logarithmic form.
\begin{smenum}
\mitemxxxx
{$5^4=625$}
{$7^{-4}=\frac{1}{2401}$}
{$\Gp{\frac{3}{4}}^3=\frac{27}{64}$}
{$13^{1/2}=\sqrt{13}$}
\end{smenum}
Write each equation in exponential form.  Use your calculator to check that the new equation you have written is true.
\begin{smenum}
\mitemxxxx
{$\log_4{64}=3$}
{$\log_3{\frac{1}{243}}=-5$}
{$\log_{1/5}{25}=-2$}
{$\log_{25}{5}=\frac{1}{2}$}
\end{smenum}

\section{Graphs of Logarithms}
Identify each graph below as that of an exponential function, a logarithmic function, or neither.
\begin{smenum}
\mitemxxxx
{\begin{ilpic}[x=.15in,y=.15in]
\draw[<->,thick] (-4,0) -- (4,0);
\draw[<->,thick] (0,-4) -- (0,4);
\clip (-3.9,-3.9) rectangle (3.9,3.9);
\draw[domain=-3.9:-.01] plot (\x,{-1/pow(2,-\x)+1});
\draw[domain=.01:3.9] plot (\x,{-pow(2,\x)+1});
\draw[dashed] (-3.9,1) -- (3.9,1);
\end{ilpic}}
{\begin{ilpic}[x=.15in,y=.15in]
\draw[<->,thick] (-4,0) -- (4,0);
\draw[<->,thick] (0,-4) -- (0,4);
\clip (-3.9,-3.9) rectangle (3.9,3.9);
\draw[domain=-3.9:-.01] plot ({1/pow(2,-\x)+1},\x);
\draw[domain=.01:3.9] plot ({pow(2,\x)+1},\x);
\draw[dashed] (1,-3.9) -- (1,3.9);
\end{ilpic}}
{\begin{ilpic}[x=.15in,y=.15in]
\draw[<->,thick] (-4,0) -- (4,0);
\draw[<->,thick] (0,-4) -- (0,4);
\clip (-3.9,-3.9) rectangle (3.9,3.9);
\draw[domain=-4.9:-.01] plot (\x+1,{1/\x+-1});
\draw[domain=.01:2.9] plot (\x+1,{1/\x+-1});
\draw[dashed] (-3.9,-1) -- (3.9,-1);
\draw[dashed] (1,-3.9) -- (1,3.9);
\end{ilpic}}
{\begin{ilpic}[x=.15in,y=.15in]
\draw[<->,thick] (-4,0) -- (4,0);
\draw[<->,thick] (0,-4) -- (0,4);
\clip (-3.9,-3.9) rectangle (3.9,3.9);
\draw[domain=-3.9:-.01] plot ({1/pow(5/3,-\x)-2},-\x);
\draw[domain=.01:3.9] plot ({pow(5/3,\x)-2},-\x);
\draw[dashed] (-2,-3.9) -- (-2,3.9);
\end{ilpic}
}
\end{smenum}
Of the two graphs above which are graphs of logarithm functions, choose the graph of each function below:
\begin{smenum}
\mitemxx
{$f(x)=\log_2(x-1)$}
{$g(x)=\log_{3/5}(x+2)$}
\end{smenum}

\section{Graph a Logarithm}
Sketch the graph of each logarithm function below in the space provided. Your graph should show the three points corresponding to outputs (not $y$ values) $-1$, $0$, and $1$, as well as the asymptote. Do your best to estimate fractions.
\begin{smenum}
\mitemxx
{$f(x)=\log_3(x)$\\\GridSquareFitEO[10]}
{$f(x)=\log_4(2-x)+2$\\\GridSquareFitEO[10]}
\mitemxx
{$f(x)=2\log_{1/2}(x+3)-1$\\\GridSquareFitEO[10]}
{$f(x)=-\log_2(3-x)+1$\\\GridSquareFitEO[10]}
\end{smenum}

\section{Domain and Range of a Logarithm Function}
For each function below, give the domain and range.
\begin{smenum}
\mitemxxx
{$h(x)=\log_{12}(x+3)+5$}
{$k(x)=\log_{1/2}(5-x)-1$}
{$u(x)=\frac{1}{3}\log_{2}(4x-12)+7$}
\end{smenum}

\section{Common Logarithm}
For each log below, rewrite in exponential form, then determine whether the sentence is true.
\begin{smenum}
\mitemxxxx
{$\log 100=2$}
{$\log 0.001=-2$}
{$\log 10000=6$}
{$\log 0.00001=-5$}
\end{smenum}
Calculate each log on your calculator. Round your answer to two decimal places if necessary. Check by rewriting in exponential form.
\begin{smenum}
\mitemxxxx
{$\log 12$}
{$\log 13,225$}
{$\log 0.03$}
{$\log 1000$}
\end{smenum}
%
%\section{Solving Exponential Equations Using Logarithms --- Basic}
%For each equation:
%\begin{clist}
%\item Use trial and error to find what two whole numbers the solution lies between.
%\item Use logarithms on your calculator to find the solution to two decimal places.
%\end{clist}
%\begin{smenum}
%\mitemxxxx
%{$2^x=25$}
%{$11\cdot 1.5^x=121$}
%{$2000\cdot 1.08^x=5000$}
%{$650\Gp{\frac{1}{2}}^x=13$}
%\end{smenum}

\section{Solving Equations Involving Logarithms --- Basic}
Solve each equation by rewriting in exponential notation.  Round your answer to two decimal places, if necessary.
\begin{smenum}
\mitemxxxx
{$\log_3 x=2$}
{$\log_8 x = \frac{4}{3}$}
{$\log_x 25 = 2$}
{$\log_3(4x+1)=4$}
\mitemxxxx
{$\log_x \frac{1}{9}=-2$}
{$\log p = 4$}
{$\log_{x+2} 9 = 2$}
{$\log x = 1.1$}
\end{smenum}

\end{document}